\section{Discussione}\label{sec:discussion}

Il passo del propagatore modifica la geometria dell'aggiornamento. Nel regime sottosmorzato, aumentare $dt$ accentua la contrazione dell'oscillazione e cambia la fase a cui viene valutata la particella (Appendice~\ref{app:dynamics}). Il confronto su più budget mostra perché la scelta di $dt$ va riferita a una configurazione e a un orizzonte di ricerca. Il valore $0.38125$ è il riferimento per lo studio di scalabilità. \rev{Nel confronto con le PSO si usa una configurazione Random-8 con casualità per coordinata, $dt=1$, ricerca locale e riavvio parziale.}

Il confronto sulla diffusione inversa individua il beneficio di distribuire le valutazioni e mostra quanto si guadagna aumentando il numero di worker. Con 512 passi RK4, master--slave e particle-parallel raggiungono accelerazioni simili rispetto alla seriale. Il passaggio a otto worker offre un ulteriore risparmio di tempo e riduce l'efficienza per worker. L'asincronia anticipa la disponibilità di alcuni attrattori, e i suoi effetti sulla qualità e sui tempi dipendono da dimensione e vicinato.

\begingroup\color{RevisionBlue}
Il confronto a tempo fissato mostra due aspetti complementari. A $D=20$ il CPSO ha il miglior tempo mediano al target, mentre a $D=100$ completa tutte le repliche entro la scadenza e si colloca dietro CSO per tempo mediano. Le FE al target indicano che, a $D=100$, CSO raggiunge la precisione richiesta con meno valutazioni. Il confronto a budget di valutazioni fissato colloca il CPSO al quarto posto in tutte e tre le suite, con i risultati migliori distribuiti su un sottoinsieme delle funzioni esterne. Le due misure distinguono la qualità della soluzione a costo di valutazione assegnato dalla rapidità con cui si raggiunge una precisione operativa.

\endgroup

La conferma dello scaling riguarda la diffusione inversa su una piattaforma con core eterogenei. \rev{Nel confronto a tempo fissato dello Studio III, tutti i metodi usano lo stesso valutatore della diffusione inversa. I tempi includono anche i costi di aggiornamento delle rispettive implementazioni.} La variazione dei passi RK4 modifica anche la discretizzazione del problema inverso e va interpretata entro la famiglia di carichi misurata.
